

\section{Produzione ed elaborazione dei dati}

\begin{center}
	\textit{Quali sono i volumi di dati in 
		gioco? Il dominio richiede risposte deterministiche real-time (bassa 
		latenza, es. Industrial/Medical) o tollera trasmissioni asincrone (es. 
		Agricoltura/Smart City)?}
\end{center}


Produzione delle informazioni locale o meno? edge IoT?

\section{Topologia e Ambiente operativo}
\begin{center}
	\textit{I dispositivi sono fissi o in mobilità 
		(es. Fleet Monitoring)? Operano in ambienti ostili (temperature 
		estreme, interferenze elettromagnetiche)?}
\end{center}
La collocazione dei sensori è strettamente legata alla specifica applicazione nel settore Agri-IoT. Nel monitoraggio delle risorse idriche o delle colture (\textit{crop monitoring}), i dispositivi vengono installati in postazioni fisse predeterminate. Al contrario, per applicazioni dedicate al controllo dei macchinari agricoli o delle attrezzature per la lavorazione del terreno, è necessario implementare soluzioni di tracciamento basate sul \textbf{fleet monitoring}, analogamente a quanto avviene nei servizi logistici.

Indipendentemente dalla configurazione, poiché tali dispositivi operano in ambienti esterni esposti alle intemperie, la loro progettazione deve privilegiare la \textbf{robustezza strutturale}. È fondamentale garantire un elevato grado di protezione contro l'infiltrazione di umidità, polvere, residui terrosi o agenti chimici, fattori che potrebbero causare danni irreversibili alla componentistica elettronica dei nodi IoT.

\section{Vincoli e consumi energetici}
\begin{center}
	\textit{I dispositivi opereranno 
		prevalentemente collegati alla rete elettrica o richiederanno 
		un'alimentazione a batteria? Quali pattern energetici (es. Event-Driven, 
		Batch Transmission, Duty-cycle) si rendono necessari?}
\end{center}

\cite{kokten2020low}

\cite{routis2023low}

\section{Sicurezza e privacy}
\begin{center}
	\textit{Qual è il livello di criticità del dato (es. dati sanitari 
		nel Medical)? Quali minacce specifiche caratterizzano questo 
		scenario?}
\end{center}

Possibili attacchi alla mole di sensori e microcontrollori può essere anche utilizzato come modo di \textit{blackmailing} da parte di gruppi terroristici o organizzazioni filo governative.
\cite{kristen2021security}



